%====================================================================
%   SECTION 2 — DATA
%====================================================================

\section{Data and pre-processing}\label{sec:data}

We use high-frequency bid--ask FX tick data from \textcite{dukascopy}. The full pipeline (data processing, back-test driver, and all figures in this thesis) is implemented in Python and released as the open-source repository \texttt{stk-mat2011} at \url{https://github.com/egil10/stk-mat2011}; implementation details, package versions, and reproduction scripts are documented there. The two legs of a pair are indexed by $i \in \{A, B\}$. The three pair constructions are AUD/USD--NZD/USD, EUR/NOK--EUR/SEK, and GBP/USD--EUR/USD, chosen to span a range of liquidity, transaction costs, and pair-dependence \parencite{krauss2017statarb}.

Rather than running one continuous back-test, we study three independent one-month sub-samples: August 2024, December 2024, and April 2025. These months were chosen because all three pairs are actively moving in them, and because they cover very different macro backdrops: summer dollar weakness, year-end positioning, and the spring 2025 tariff episode. Each trading day inside a month is its own self-contained back-test. Every model and threshold introduced later (the rolling cointegration in Section~\ref{sec:spread}, the $z$-score in Section~\ref{sec:zscore}, and the regime model in Section~\ref{sec:msarspread}) is re-fitted from scratch at the start of every session, so any in-sample look-ahead is bounded by one day.

\paragraph{From bid and ask to mid and cost.}
Bid and ask streams are sorted by timestamp and aligned by a backward as-of merge. The bid and ask are used in two roles, never both at once. The \emph{mid-price} (the average of bid and ask) is the input to every model in the rest of the thesis; it is the price the strategy sees when it decides whether to trade. The \emph{half-spread} (half the bid--ask gap, in basis points) is not fed to any model; it is used in the back-test of Section~\ref{sec:walkforward} as the per-leg cost paid whenever the strategy crosses the market. Writing $P_t^{\mathrm{bid},\,i}$ and $P_t^{\mathrm{ask},\,i}$ for the prevailing bid and ask of leg $i$ at time $t$,
\begin{equation}\label{eq:midprice}
\begin{aligned}
    P_t^{\mathrm{mid},\,i} \;&=\; \tfrac{1}{2}\bigl(P_t^{\mathrm{bid},\,i} + P_t^{\mathrm{ask},\,i}\bigr),\\[3pt]
    \mathrm{HS}_t^{\,i} \;&=\; \tfrac{1}{2}\,\frac{P_t^{\mathrm{ask},\,i} - P_t^{\mathrm{bid},\,i}}{P_t^{\mathrm{mid},\,i}}\,\times\, 10^{4}.
\end{aligned}
\end{equation}
$\mathrm{HS}_t^{\,i}$ is the cost of crossing the spread on one leg \parencite{roll1984spread}. A pairs trade enters and exits both legs at the same time, so the combined per-flip cost is $\mathrm{HS}_t^{\,A} + \hat\beta\,\mathrm{HS}_t^{\,B}$; whether the rolling rule of Section~\ref{sec:zscore} can clear that cost is the central question of the empirical part of the thesis.

\paragraph{Pre-averaging on number of ticks.}\label{sec:preaveraging}
The raw mid-price is noisy on a tick-by-tick basis: bid--ask bounce, discreteness, and asynchronous quote updates all push the observed price around even when the underlying ``efficient'' price is unchanged \parencite{hasbrouck2007microstructure}. We reduce this noise by pre-averaging in tick time, following \textcite{jacod2017}: we group consecutive ticks into blocks of $L$ ticks each and replace every block by the average mid-price inside it. If $P_{\tau}^{\mathrm{mid},\,i}$ is the raw tick-level mid-price of leg $i$ at tick $\tau$, the pre-averaged mid-price at block index $t$ is
\begin{equation}\label{eq:preaverage}
    \overline{P}_t^{\,i} \;=\; \frac{1}{L}\sum_{\tau=(t-1)L+1}^{tL} P_{\tau}^{\mathrm{mid},\,i},
    \qquad t = 1, 2, \ldots
\end{equation}
The same block size $L$ is used for both legs, so the pre-averaged legs stay synchronised. Pre-averaging in tick time is preferred to fixed-clock sampling because it gives an observation index that is more even in microstructure activity and lower in noise \parencite{jacod2017}. From here on, every quantity built from the mid-price (log-returns, cointegration spread, $z$-score, HMM observations) is built from $\overline{P}_t^{\,i}$, with $t$ indexing blocks rather than ticks. To keep the notation light, we write $P_t^{\mathrm{mid},\,i} \equiv \overline{P}_t^{\,i}$ in Sections~\ref{sec:pairsbaseline}--\ref{sec:result}.

